\documentclass[11pt]{article}
% ---------------------------------------------------------------
% R4.02 Qualite de developpement -- preambule commun des sujets de TP et des corrections
% Article clair, fond blanc, accents de couleur coherents avec les slides
% Charte reprise de R3.07 (meme palette IUT que les diapositives)
% Utilisation :
%   \documentclass[11pt]{article}
%   % ---------------------------------------------------------------
% R4.02 Qualite de developpement -- preambule commun des sujets de TP et des corrections
% Article clair, fond blanc, accents de couleur coherents avec les slides
% Charte reprise de R3.07 (meme palette IUT que les diapositives)
% Utilisation :
%   \documentclass[11pt]{article}
%   % ---------------------------------------------------------------
% R4.02 Qualite de developpement -- preambule commun des sujets de TP et des corrections
% Article clair, fond blanc, accents de couleur coherents avec les slides
% Charte reprise de R3.07 (meme palette IUT que les diapositives)
% Utilisation :
%   \documentclass[11pt]{article}
%   \input{../../latex/common/r402-tp-preamble.tex}
%   \begin{document}
%   \rtpheader{1}{Titre du TP}
% ---------------------------------------------------------------
\usepackage[T1]{fontenc}
\usepackage[utf8]{inputenc}
\usepackage[french]{babel}
\usepackage[scale=0.92]{tgheros}
\usepackage{tgpagella}
\usepackage[a4paper, margin=2.2cm, headheight=14pt]{geometry}
\usepackage{graphicx}
\usepackage{xcolor}
\usepackage{listings}
\usepackage[skins,listings,breakable]{tcolorbox}
\usepackage{enumitem}
\usepackage{booktabs}
\usepackage{array}
\usepackage{multirow}
\usepackage{tikz}
\usetikzlibrary{shapes, shapes.geometric, arrows, arrows.meta, positioning, calc}
\usepackage[colorlinks=true, linkcolor=qaaccentdark, urlcolor=qaaccentdark]{hyperref}
\usepackage{fancyhdr}
\usepackage{lastpage}
\usepackage{titlesec}
\usepackage{etoolbox}
\usepackage{eurosym}

% ---------------------------------------------------------------
% MENTION D'ELABORATION (discrete : gris clair, petit corps)
% ---------------------------------------------------------------
\definecolor{qaaimention}{HTML}{A3AAB5}
\newcommand{\raimention}{{\scriptsize\color{qaaimention}Supports \'elabor\'es avec l'aide de Claude Code Opus 5.5}}

% ---------------------------------------------------------------
% PALETTE -- version imprimable des couleurs des diapositives
% ---------------------------------------------------------------
\definecolor{qaaccentdark}{HTML}{7A1680}   % violet du logo IUT, fonce pour rester lisible sur blanc
\definecolor{qaaccent2dark}{HTML}{9B4B94}  % violet clair du logo IUT (accent secondaire)
\definecolor{qapanellight}{HTML}{F4F6F8}   % fond de bloc tres clair
\definecolor{qalinelight}{HTML}{D8DEE4}    % filets
\definecolor{qatextdark}{HTML}{1A1D23}     % texte principal
\definecolor{qamuteddark}{HTML}{5B6472}    % texte secondaire
\definecolor{qacodekwl}{HTML}{7A1680}      % code : mots-cles
\definecolor{qacodestrl}{HTML}{2E7D32}     % code : chaines
\definecolor{qacodecoml}{HTML}{6B7280}     % code : commentaires
\definecolor{qacodeannl}{HTML}{B26A00}     % code : annotations / javadoc

\color{qatextdark}

\setlength{\parskip}{4pt}
\setlength{\parindent}{0pt}

% ---------------------------------------------------------------
% TITRES DE SECTION -- texte blanc sur bande de couleur pleine largeur
% (gere aussi les variantes etoilees \section*)
% ---------------------------------------------------------------
\newcommand{\qasectionband}[2]{%
  \noindent\colorbox{#1}{%
    \parbox[c]{\dimexpr\linewidth-2\fboxsep\relax}{%
      \color{white}#2%
    }%
  }%
}
\newcommand{\secformat}[1]{\qasectionband{qaaccentdark}{\bfseries\Large\thesection\quad #1}}
\newcommand{\secformatnn}[1]{\qasectionband{qaaccentdark}{\bfseries\Large #1}}
\newcommand{\subsecformat}[1]{\qasectionband{qaaccent2dark}{\bfseries\large\thesubsection\quad #1}}
\newcommand{\subsecformatnn}[1]{\qasectionband{qaaccent2dark}{\bfseries\large #1}}

\titleformat{\section}{\normalfont}{}{0pt}{\secformat}
\titleformat{name=\section,numberless}{\normalfont}{}{0pt}{\secformatnn}
\titleformat{\subsection}{\normalfont}{}{0pt}{\subsecformat}
\titleformat{name=\subsection,numberless}{\normalfont}{}{0pt}{\subsecformatnn}
\titleformat{\subsubsection}{\normalfont\bfseries\color{qaaccentdark}}{\thesubsubsection}{0.6em}{}
\titlespacing*{\section}{0pt}{14pt}{8pt}
\titlespacing*{\subsection}{0pt}{10pt}{6pt}
\titlespacing*{\subsubsection}{0pt}{8pt}{4pt}

% ---------------------------------------------------------------
% LISTINGS -- theme clair (lstlisting reste utilisable tel quel)
% \lstcode : police + couleur de base, a reutiliser dans les surcharges
% ---------------------------------------------------------------
\newcommand{\lstcode}{\ttfamily\color{qatextdark}}

\lstdefinelanguage{TypeScript}{
  keywords={typeof, new, true, false, catch, function, return, null, switch, var,
    if, in, while, do, else, case, break, const, let, interface, class, extends,
    implements, enum, export, import, this, super, constructor, public, private,
    protected, static, async, await, throw, throws, try, finally, void},
  ndkeywords={boolean, any, number, string, int, long, double},
  sensitive=true,
  comment=[l]{//},
  morecomment=[s]{/*}{*/},
  morestring=[b]',
  morestring=[b]"
}

\lstdefinestyle{r402code}{
  basicstyle=\lstcode\footnotesize,
  keywordstyle=\color{qacodekwl}\bfseries,
  ndkeywordstyle=\color{qaaccent2dark}\bfseries,
  stringstyle=\color{qacodestrl},
  commentstyle=\color{qacodecoml}\itshape,
  morecomment=[s][\color{qacodeannl}]{/**}{*/},
  backgroundcolor=\color{qapanellight},
  rulecolor=\color{qalinelight},
  frame=single,
  framesep=4pt,
  xleftmargin=6pt,
  xrightmargin=6pt,
  numbers=left,
  numberstyle=\tiny\color{qamuteddark},
  numbersep=8pt,
  keepspaces=true,
  columns=fullflexible,
  showstringspaces=false,
  showspaces=false,
  showtabs=false,
  breaklines=true,
  breakatwhitespace=false,
  captionpos=b,
  tabsize=2,
  escapeinside={(*@}{@*)},
  % accents UTF-8 dans le code (listings ne gere pas l'UTF-8 multi-octets)
  literate={é}{{\'e}}1 {è}{{\`e}}1 {ê}{{\^e}}1 {à}{{\`a}}1 {â}{{\^a}}1
           {ù}{{\`u}}1 {û}{{\^u}}1 {ô}{{\^o}}1 {î}{{\^i}}1 {ç}{{\c c}}1
           {É}{{\'E}}1 {À}{{\`A}}1 {€}{{\euro{}}}1 {…}{{\ldots}}1,
}
\lstset{style=r402code}

% ---------------------------------------------------------------
% BOITES : exercice, code, information
% ---------------------------------------------------------------
\tcbset{
  exobox/.style={colback=qapanellight, colframe=qaaccentdark, boxrule=0.6pt,
    arc=2pt, left=8pt, right=8pt, top=6pt, bottom=6pt,
    coltitle=white, colbacktitle=qaaccentdark, fonttitle=\bfseries,
    boxsep=2pt, breakable},
  infobox/.style={colback=white, colframe=qaaccent2dark, boxrule=0.6pt,
    leftrule=3pt, arc=0pt, left=8pt, right=8pt, top=5pt, bottom=5pt,
    coltitle=qaaccent2dark, colbacktitle=white, fonttitle=\bfseries,
    breakable},
  codebox/.style={colback=qatextdark, colframe=qalinelight, boxrule=0.4pt,
    arc=2pt, left=8pt, right=8pt, top=5pt, bottom=5pt, listing only,
    listing options={basicstyle=\ttfamily\small\color{white},
      breaklines=true, showstringspaces=false, columns=fullflexible,
      keepspaces=true, frame=none, numbers=none, backgroundcolor={}}, breakable}
}
\newtcolorbox[auto counter]{exobox}[2][]{exobox,
  title={Exercice~\thetcbcounter\ifstrempty{#2}{}{ --- #2}}, #1}
\newtcolorbox{infobox}[1][\`A retenir]{infobox, title={#1}}
\newtcblisting{codebox}{codebox}

% ---------------------------------------------------------------
% EN-TETE DE PAGE (fil d'Ariane + pagination)
% ---------------------------------------------------------------
\newcommand{\qaheadtitlevalue}{}
\newcommand{\qaheadtitle}[1]{\renewcommand{\qaheadtitlevalue}{#1}}

\pagestyle{fancy}
\fancyhf{}
\renewcommand{\headrulewidth}{0.4pt}
\fancyhead[L]{\small\color{qamuteddark}R4.02 --- \qaheadtitlevalue}
\fancyhead[R]{\small\color{qamuteddark}\thepage\,/\,\pageref{LastPage}}
\fancyfoot[C]{\footnotesize\color{qamuteddark}IUT de Vannes --- Universit\'e Bretagne Sud}
\fancypagestyle{plain}{\fancyhf{}\renewcommand{\headrulewidth}{0pt}%
  \fancyfoot[C]{\footnotesize\color{qamuteddark}IUT de Vannes --- Universit\'e Bretagne Sud}}

% ---------------------------------------------------------------
% BANDEAU DE TITRE (pas une page dediee, pour rester compact)
% \rtpheader{numero ou intitule court}{titre du TP}
% ---------------------------------------------------------------
\newcommand{\ryear}{2026/2027}
\newcommand{\rtpheader}[2]{%
  \thispagestyle{plain}%
  \begin{minipage}[t]{0.68\textwidth}
    {\small\color{qamuteddark}
      Universit\'e Bretagne Sud\\
      IUT de Vannes\\
      D\'epartement Informatique --- \ryear}
  \end{minipage}%
  \hfill
  \begin{minipage}[t]{0.28\textwidth}
    \raggedleft
    \includegraphics[width=2.6cm]{iut_logo.jpg}
  \end{minipage}

  \vspace{6mm}
  \begin{center}
    {\large\bfseries\color{qatextdark} BUT 2 --- R4.02. Qualit\'e de d\'eveloppement}\\[2mm]
    {\normalsize\color{qamuteddark} TP #1}\\[1mm]
    {\Large\bfseries\color{qaaccentdark} #2}\\[2mm]
    \raimention
  \end{center}
  \vspace{4mm}
  \qaheadtitle{#2}
}

%   \begin{document}
%   \rtpheader{1}{Titre du TP}
% ---------------------------------------------------------------
\usepackage[T1]{fontenc}
\usepackage[utf8]{inputenc}
\usepackage[french]{babel}
\usepackage[scale=0.92]{tgheros}
\usepackage{tgpagella}
\usepackage[a4paper, margin=2.2cm, headheight=14pt]{geometry}
\usepackage{graphicx}
\usepackage{xcolor}
\usepackage{listings}
\usepackage[skins,listings,breakable]{tcolorbox}
\usepackage{enumitem}
\usepackage{booktabs}
\usepackage{array}
\usepackage{multirow}
\usepackage{tikz}
\usetikzlibrary{shapes, shapes.geometric, arrows, arrows.meta, positioning, calc}
\usepackage[colorlinks=true, linkcolor=qaaccentdark, urlcolor=qaaccentdark]{hyperref}
\usepackage{fancyhdr}
\usepackage{lastpage}
\usepackage{titlesec}
\usepackage{etoolbox}
\usepackage{eurosym}

% ---------------------------------------------------------------
% MENTION D'ELABORATION (discrete : gris clair, petit corps)
% ---------------------------------------------------------------
\definecolor{qaaimention}{HTML}{A3AAB5}
\newcommand{\raimention}{{\scriptsize\color{qaaimention}Supports \'elabor\'es avec l'aide de Claude Code Opus 5.5}}

% ---------------------------------------------------------------
% PALETTE -- version imprimable des couleurs des diapositives
% ---------------------------------------------------------------
\definecolor{qaaccentdark}{HTML}{7A1680}   % violet du logo IUT, fonce pour rester lisible sur blanc
\definecolor{qaaccent2dark}{HTML}{9B4B94}  % violet clair du logo IUT (accent secondaire)
\definecolor{qapanellight}{HTML}{F4F6F8}   % fond de bloc tres clair
\definecolor{qalinelight}{HTML}{D8DEE4}    % filets
\definecolor{qatextdark}{HTML}{1A1D23}     % texte principal
\definecolor{qamuteddark}{HTML}{5B6472}    % texte secondaire
\definecolor{qacodekwl}{HTML}{7A1680}      % code : mots-cles
\definecolor{qacodestrl}{HTML}{2E7D32}     % code : chaines
\definecolor{qacodecoml}{HTML}{6B7280}     % code : commentaires
\definecolor{qacodeannl}{HTML}{B26A00}     % code : annotations / javadoc

\color{qatextdark}

\setlength{\parskip}{4pt}
\setlength{\parindent}{0pt}

% ---------------------------------------------------------------
% TITRES DE SECTION -- texte blanc sur bande de couleur pleine largeur
% (gere aussi les variantes etoilees \section*)
% ---------------------------------------------------------------
\newcommand{\qasectionband}[2]{%
  \noindent\colorbox{#1}{%
    \parbox[c]{\dimexpr\linewidth-2\fboxsep\relax}{%
      \color{white}#2%
    }%
  }%
}
\newcommand{\secformat}[1]{\qasectionband{qaaccentdark}{\bfseries\Large\thesection\quad #1}}
\newcommand{\secformatnn}[1]{\qasectionband{qaaccentdark}{\bfseries\Large #1}}
\newcommand{\subsecformat}[1]{\qasectionband{qaaccent2dark}{\bfseries\large\thesubsection\quad #1}}
\newcommand{\subsecformatnn}[1]{\qasectionband{qaaccent2dark}{\bfseries\large #1}}

\titleformat{\section}{\normalfont}{}{0pt}{\secformat}
\titleformat{name=\section,numberless}{\normalfont}{}{0pt}{\secformatnn}
\titleformat{\subsection}{\normalfont}{}{0pt}{\subsecformat}
\titleformat{name=\subsection,numberless}{\normalfont}{}{0pt}{\subsecformatnn}
\titleformat{\subsubsection}{\normalfont\bfseries\color{qaaccentdark}}{\thesubsubsection}{0.6em}{}
\titlespacing*{\section}{0pt}{14pt}{8pt}
\titlespacing*{\subsection}{0pt}{10pt}{6pt}
\titlespacing*{\subsubsection}{0pt}{8pt}{4pt}

% ---------------------------------------------------------------
% LISTINGS -- theme clair (lstlisting reste utilisable tel quel)
% \lstcode : police + couleur de base, a reutiliser dans les surcharges
% ---------------------------------------------------------------
\newcommand{\lstcode}{\ttfamily\color{qatextdark}}

\lstdefinelanguage{TypeScript}{
  keywords={typeof, new, true, false, catch, function, return, null, switch, var,
    if, in, while, do, else, case, break, const, let, interface, class, extends,
    implements, enum, export, import, this, super, constructor, public, private,
    protected, static, async, await, throw, throws, try, finally, void},
  ndkeywords={boolean, any, number, string, int, long, double},
  sensitive=true,
  comment=[l]{//},
  morecomment=[s]{/*}{*/},
  morestring=[b]',
  morestring=[b]"
}

\lstdefinestyle{r402code}{
  basicstyle=\lstcode\footnotesize,
  keywordstyle=\color{qacodekwl}\bfseries,
  ndkeywordstyle=\color{qaaccent2dark}\bfseries,
  stringstyle=\color{qacodestrl},
  commentstyle=\color{qacodecoml}\itshape,
  morecomment=[s][\color{qacodeannl}]{/**}{*/},
  backgroundcolor=\color{qapanellight},
  rulecolor=\color{qalinelight},
  frame=single,
  framesep=4pt,
  xleftmargin=6pt,
  xrightmargin=6pt,
  numbers=left,
  numberstyle=\tiny\color{qamuteddark},
  numbersep=8pt,
  keepspaces=true,
  columns=fullflexible,
  showstringspaces=false,
  showspaces=false,
  showtabs=false,
  breaklines=true,
  breakatwhitespace=false,
  captionpos=b,
  tabsize=2,
  escapeinside={(*@}{@*)},
  % accents UTF-8 dans le code (listings ne gere pas l'UTF-8 multi-octets)
  literate={é}{{\'e}}1 {è}{{\`e}}1 {ê}{{\^e}}1 {à}{{\`a}}1 {â}{{\^a}}1
           {ù}{{\`u}}1 {û}{{\^u}}1 {ô}{{\^o}}1 {î}{{\^i}}1 {ç}{{\c c}}1
           {É}{{\'E}}1 {À}{{\`A}}1 {€}{{\euro{}}}1 {…}{{\ldots}}1,
}
\lstset{style=r402code}

% ---------------------------------------------------------------
% BOITES : exercice, code, information
% ---------------------------------------------------------------
\tcbset{
  exobox/.style={colback=qapanellight, colframe=qaaccentdark, boxrule=0.6pt,
    arc=2pt, left=8pt, right=8pt, top=6pt, bottom=6pt,
    coltitle=white, colbacktitle=qaaccentdark, fonttitle=\bfseries,
    boxsep=2pt, breakable},
  infobox/.style={colback=white, colframe=qaaccent2dark, boxrule=0.6pt,
    leftrule=3pt, arc=0pt, left=8pt, right=8pt, top=5pt, bottom=5pt,
    coltitle=qaaccent2dark, colbacktitle=white, fonttitle=\bfseries,
    breakable},
  codebox/.style={colback=qatextdark, colframe=qalinelight, boxrule=0.4pt,
    arc=2pt, left=8pt, right=8pt, top=5pt, bottom=5pt, listing only,
    listing options={basicstyle=\ttfamily\small\color{white},
      breaklines=true, showstringspaces=false, columns=fullflexible,
      keepspaces=true, frame=none, numbers=none, backgroundcolor={}}, breakable}
}
\newtcolorbox[auto counter]{exobox}[2][]{exobox,
  title={Exercice~\thetcbcounter\ifstrempty{#2}{}{ --- #2}}, #1}
\newtcolorbox{infobox}[1][\`A retenir]{infobox, title={#1}}
\newtcblisting{codebox}{codebox}

% ---------------------------------------------------------------
% EN-TETE DE PAGE (fil d'Ariane + pagination)
% ---------------------------------------------------------------
\newcommand{\qaheadtitlevalue}{}
\newcommand{\qaheadtitle}[1]{\renewcommand{\qaheadtitlevalue}{#1}}

\pagestyle{fancy}
\fancyhf{}
\renewcommand{\headrulewidth}{0.4pt}
\fancyhead[L]{\small\color{qamuteddark}R4.02 --- \qaheadtitlevalue}
\fancyhead[R]{\small\color{qamuteddark}\thepage\,/\,\pageref{LastPage}}
\fancyfoot[C]{\footnotesize\color{qamuteddark}IUT de Vannes --- Universit\'e Bretagne Sud}
\fancypagestyle{plain}{\fancyhf{}\renewcommand{\headrulewidth}{0pt}%
  \fancyfoot[C]{\footnotesize\color{qamuteddark}IUT de Vannes --- Universit\'e Bretagne Sud}}

% ---------------------------------------------------------------
% BANDEAU DE TITRE (pas une page dediee, pour rester compact)
% \rtpheader{numero ou intitule court}{titre du TP}
% ---------------------------------------------------------------
\newcommand{\ryear}{2026/2027}
\newcommand{\rtpheader}[2]{%
  \thispagestyle{plain}%
  \begin{minipage}[t]{0.68\textwidth}
    {\small\color{qamuteddark}
      Universit\'e Bretagne Sud\\
      IUT de Vannes\\
      D\'epartement Informatique --- \ryear}
  \end{minipage}%
  \hfill
  \begin{minipage}[t]{0.28\textwidth}
    \raggedleft
    \includegraphics[width=2.6cm]{iut_logo.jpg}
  \end{minipage}

  \vspace{6mm}
  \begin{center}
    {\large\bfseries\color{qatextdark} BUT 2 --- R4.02. Qualit\'e de d\'eveloppement}\\[2mm]
    {\normalsize\color{qamuteddark} TP #1}\\[1mm]
    {\Large\bfseries\color{qaaccentdark} #2}\\[2mm]
    \raimention
  \end{center}
  \vspace{4mm}
  \qaheadtitle{#2}
}

%   \begin{document}
%   \rtpheader{1}{Titre du TP}
% ---------------------------------------------------------------
\usepackage[T1]{fontenc}
\usepackage[utf8]{inputenc}
\usepackage[french]{babel}
\usepackage[scale=0.92]{tgheros}
\usepackage{tgpagella}
\usepackage[a4paper, margin=2.2cm, headheight=14pt]{geometry}
\usepackage{graphicx}
\usepackage{xcolor}
\usepackage{listings}
\usepackage[skins,listings,breakable]{tcolorbox}
\usepackage{enumitem}
\usepackage{booktabs}
\usepackage{array}
\usepackage{multirow}
\usepackage{tikz}
\usetikzlibrary{shapes, shapes.geometric, arrows, arrows.meta, positioning, calc}
\usepackage[colorlinks=true, linkcolor=qaaccentdark, urlcolor=qaaccentdark]{hyperref}
\usepackage{fancyhdr}
\usepackage{lastpage}
\usepackage{titlesec}
\usepackage{etoolbox}
\usepackage{eurosym}

% ---------------------------------------------------------------
% MENTION D'ELABORATION (discrete : gris clair, petit corps)
% ---------------------------------------------------------------
\definecolor{qaaimention}{HTML}{A3AAB5}
\newcommand{\raimention}{{\scriptsize\color{qaaimention}Supports \'elabor\'es avec l'aide de Claude Code Opus 5.5}}

% ---------------------------------------------------------------
% PALETTE -- version imprimable des couleurs des diapositives
% ---------------------------------------------------------------
\definecolor{qaaccentdark}{HTML}{7A1680}   % violet du logo IUT, fonce pour rester lisible sur blanc
\definecolor{qaaccent2dark}{HTML}{9B4B94}  % violet clair du logo IUT (accent secondaire)
\definecolor{qapanellight}{HTML}{F4F6F8}   % fond de bloc tres clair
\definecolor{qalinelight}{HTML}{D8DEE4}    % filets
\definecolor{qatextdark}{HTML}{1A1D23}     % texte principal
\definecolor{qamuteddark}{HTML}{5B6472}    % texte secondaire
\definecolor{qacodekwl}{HTML}{7A1680}      % code : mots-cles
\definecolor{qacodestrl}{HTML}{2E7D32}     % code : chaines
\definecolor{qacodecoml}{HTML}{6B7280}     % code : commentaires
\definecolor{qacodeannl}{HTML}{B26A00}     % code : annotations / javadoc

\color{qatextdark}

\setlength{\parskip}{4pt}
\setlength{\parindent}{0pt}

% ---------------------------------------------------------------
% TITRES DE SECTION -- texte blanc sur bande de couleur pleine largeur
% (gere aussi les variantes etoilees \section*)
% ---------------------------------------------------------------
\newcommand{\qasectionband}[2]{%
  \noindent\colorbox{#1}{%
    \parbox[c]{\dimexpr\linewidth-2\fboxsep\relax}{%
      \color{white}#2%
    }%
  }%
}
\newcommand{\secformat}[1]{\qasectionband{qaaccentdark}{\bfseries\Large\thesection\quad #1}}
\newcommand{\secformatnn}[1]{\qasectionband{qaaccentdark}{\bfseries\Large #1}}
\newcommand{\subsecformat}[1]{\qasectionband{qaaccent2dark}{\bfseries\large\thesubsection\quad #1}}
\newcommand{\subsecformatnn}[1]{\qasectionband{qaaccent2dark}{\bfseries\large #1}}

\titleformat{\section}{\normalfont}{}{0pt}{\secformat}
\titleformat{name=\section,numberless}{\normalfont}{}{0pt}{\secformatnn}
\titleformat{\subsection}{\normalfont}{}{0pt}{\subsecformat}
\titleformat{name=\subsection,numberless}{\normalfont}{}{0pt}{\subsecformatnn}
\titleformat{\subsubsection}{\normalfont\bfseries\color{qaaccentdark}}{\thesubsubsection}{0.6em}{}
\titlespacing*{\section}{0pt}{14pt}{8pt}
\titlespacing*{\subsection}{0pt}{10pt}{6pt}
\titlespacing*{\subsubsection}{0pt}{8pt}{4pt}

% ---------------------------------------------------------------
% LISTINGS -- theme clair (lstlisting reste utilisable tel quel)
% \lstcode : police + couleur de base, a reutiliser dans les surcharges
% ---------------------------------------------------------------
\newcommand{\lstcode}{\ttfamily\color{qatextdark}}

\lstdefinelanguage{TypeScript}{
  keywords={typeof, new, true, false, catch, function, return, null, switch, var,
    if, in, while, do, else, case, break, const, let, interface, class, extends,
    implements, enum, export, import, this, super, constructor, public, private,
    protected, static, async, await, throw, throws, try, finally, void},
  ndkeywords={boolean, any, number, string, int, long, double},
  sensitive=true,
  comment=[l]{//},
  morecomment=[s]{/*}{*/},
  morestring=[b]',
  morestring=[b]"
}

\lstdefinestyle{r402code}{
  basicstyle=\lstcode\footnotesize,
  keywordstyle=\color{qacodekwl}\bfseries,
  ndkeywordstyle=\color{qaaccent2dark}\bfseries,
  stringstyle=\color{qacodestrl},
  commentstyle=\color{qacodecoml}\itshape,
  morecomment=[s][\color{qacodeannl}]{/**}{*/},
  backgroundcolor=\color{qapanellight},
  rulecolor=\color{qalinelight},
  frame=single,
  framesep=4pt,
  xleftmargin=6pt,
  xrightmargin=6pt,
  numbers=left,
  numberstyle=\tiny\color{qamuteddark},
  numbersep=8pt,
  keepspaces=true,
  columns=fullflexible,
  showstringspaces=false,
  showspaces=false,
  showtabs=false,
  breaklines=true,
  breakatwhitespace=false,
  captionpos=b,
  tabsize=2,
  escapeinside={(*@}{@*)},
  % accents UTF-8 dans le code (listings ne gere pas l'UTF-8 multi-octets)
  literate={é}{{\'e}}1 {è}{{\`e}}1 {ê}{{\^e}}1 {à}{{\`a}}1 {â}{{\^a}}1
           {ù}{{\`u}}1 {û}{{\^u}}1 {ô}{{\^o}}1 {î}{{\^i}}1 {ç}{{\c c}}1
           {É}{{\'E}}1 {À}{{\`A}}1 {€}{{\euro{}}}1 {…}{{\ldots}}1,
}
\lstset{style=r402code}

% ---------------------------------------------------------------
% BOITES : exercice, code, information
% ---------------------------------------------------------------
\tcbset{
  exobox/.style={colback=qapanellight, colframe=qaaccentdark, boxrule=0.6pt,
    arc=2pt, left=8pt, right=8pt, top=6pt, bottom=6pt,
    coltitle=white, colbacktitle=qaaccentdark, fonttitle=\bfseries,
    boxsep=2pt, breakable},
  infobox/.style={colback=white, colframe=qaaccent2dark, boxrule=0.6pt,
    leftrule=3pt, arc=0pt, left=8pt, right=8pt, top=5pt, bottom=5pt,
    coltitle=qaaccent2dark, colbacktitle=white, fonttitle=\bfseries,
    breakable},
  codebox/.style={colback=qatextdark, colframe=qalinelight, boxrule=0.4pt,
    arc=2pt, left=8pt, right=8pt, top=5pt, bottom=5pt, listing only,
    listing options={basicstyle=\ttfamily\small\color{white},
      breaklines=true, showstringspaces=false, columns=fullflexible,
      keepspaces=true, frame=none, numbers=none, backgroundcolor={}}, breakable}
}
\newtcolorbox[auto counter]{exobox}[2][]{exobox,
  title={Exercice~\thetcbcounter\ifstrempty{#2}{}{ --- #2}}, #1}
\newtcolorbox{infobox}[1][\`A retenir]{infobox, title={#1}}
\newtcblisting{codebox}{codebox}

% ---------------------------------------------------------------
% EN-TETE DE PAGE (fil d'Ariane + pagination)
% ---------------------------------------------------------------
\newcommand{\qaheadtitlevalue}{}
\newcommand{\qaheadtitle}[1]{\renewcommand{\qaheadtitlevalue}{#1}}

\pagestyle{fancy}
\fancyhf{}
\renewcommand{\headrulewidth}{0.4pt}
\fancyhead[L]{\small\color{qamuteddark}R4.02 --- \qaheadtitlevalue}
\fancyhead[R]{\small\color{qamuteddark}\thepage\,/\,\pageref{LastPage}}
\fancyfoot[C]{\footnotesize\color{qamuteddark}IUT de Vannes --- Universit\'e Bretagne Sud}
\fancypagestyle{plain}{\fancyhf{}\renewcommand{\headrulewidth}{0pt}%
  \fancyfoot[C]{\footnotesize\color{qamuteddark}IUT de Vannes --- Universit\'e Bretagne Sud}}

% ---------------------------------------------------------------
% BANDEAU DE TITRE (pas une page dediee, pour rester compact)
% \rtpheader{numero ou intitule court}{titre du TP}
% ---------------------------------------------------------------
\newcommand{\ryear}{2026/2027}
\newcommand{\rtpheader}[2]{%
  \thispagestyle{plain}%
  \begin{minipage}[t]{0.68\textwidth}
    {\small\color{qamuteddark}
      Universit\'e Bretagne Sud\\
      IUT de Vannes\\
      D\'epartement Informatique --- \ryear}
  \end{minipage}%
  \hfill
  \begin{minipage}[t]{0.28\textwidth}
    \raggedleft
    \includegraphics[width=2.6cm]{iut_logo.jpg}
  \end{minipage}

  \vspace{6mm}
  \begin{center}
    {\large\bfseries\color{qatextdark} BUT 2 --- R4.02. Qualit\'e de d\'eveloppement}\\[2mm]
    {\normalsize\color{qamuteddark} TP #1}\\[1mm]
    {\Large\bfseries\color{qaaccentdark} #2}\\[2mm]
    \raimention
  \end{center}
  \vspace{4mm}
  \qaheadtitle{#2}
}

\graphicspath{{../../latex/common/assets/}{./}}

% A ajuster avant diffusion
\newcommand{\rdatelimite}{la veille de la séance du TP5, à 23\,h\,59}

\begin{document}

\rtpheader{4}{Tests d'intégration avec une base de données}

\begin{infobox}[Objectifs]
\begin{itemize}[leftmargin=1.2em, itemsep=1pt]
  \item Tester les repositories Spring Data sur un vrai PostgreSQL, démarré par Testcontainers.
  \item Constater qu'une base en mémoire (H2) ne se comporte pas comme PostgreSQL.
  \item Tester l'API REST de bout en bout côté serveur (contrôleur, service, base) avec MockMvc.
  \item Garantir l'isolation des tests et vérifier qu'une transaction ne laisse pas de données incohérentes.
\end{itemize}
\end{infobox}

\begin{infobox}[Prérequis : Docker]
Testcontainers a besoin de Docker. Vérifiez avant tout que \texttt{docker ps} fonctionne sur votre poste, puis téléchargez l'image une fois pour toutes : \texttt{docker pull postgres:18-alpine}.
\end{infobox}

Continuez dans le dépôt de votre binôme, à partir de votre rendu du TP3. Le projet contient déjà tout ce qu'il faut pour ce TP : les dépendances de test (Testcontainers, MockMvc), les migrations Flyway qui créent le schéma de la base (dossier \texttt{src/main/resources/db/migration}), les contrôleurs REST (paquetage \texttt{web}) et les repositories Spring Data (paquetage \texttt{persistance}).

% ================================================================
\section{Mise en place (10 min)}

\begin{exobox}{Un premier test avec un vrai PostgreSQL}
\begin{enumerate}[leftmargin=1.4em]
  \item Créez \texttt{src/test/java/fr/iutvannes/banque/persistance/CompteRepositoryTest.java} en vous inspirant du cours : \texttt{@DataJpaTest}, \texttt{@AutoConfigureTestDatabase(replace = Replace.NONE)}, \texttt{@Testcontainers}, et un conteneur \texttt{static} annoté \texttt{@Container @ServiceConnection}.
  \item Écrivez un test qui enregistre un compte puis le retrouve avec \texttt{findByNumero}. Lancez-le et observez \texttt{docker ps} pendant l'exécution.
  \item Combien de temps prend le test ? Combien de temps prend le démarrage du conteneur ? Notez-le dans \texttt{docs/tp4.md}.
  \item Vos tests ne doivent \textbf{jamais} utiliser la base de développement (\texttt{docker compose}). Arrêtez-la (\texttt{docker compose down}) et relancez le test : passe-t-il toujours ? Pourquoi ?
\end{enumerate}
\end{exobox}

\begin{infobox}[Partager le conteneur entre plusieurs classes de test]
Un conteneur \texttt{static} est démarré une fois \textbf{par classe}. Pour le partager entre toutes les classes, déclarez-le une seule fois comme bean de test, et importez cette configuration dans chaque classe avec \texttt{@Import(ConteneursDeTest.class)} :
\begin{lstlisting}[language=Java]
@TestConfiguration(proxyBeanMethods = false)
public class ConteneursDeTest {
    @Bean
    @ServiceConnection
    PostgreSQLContainer postgres() {
        return new PostgreSQLContainer("postgres:18-alpine");
    }
}
\end{lstlisting}
\end{infobox}

% ================================================================
\section{Tester les repositories (25 min)}

Le projet contient maintenant un \texttt{OperationRepository} avec trois méthodes, utilisées pour produire le relevé mensuel d'un compte :

\begin{lstlisting}[language=Java]
public interface OperationRepository extends Repository<Operation, Long> {

    Operation save(Operation operation);

    // Operations d'un compte entre deux instants, de la plus recente a la plus ancienne
    List<Operation> findByCompteNumeroAndDateBetweenOrderByDateDesc(
            String numero, LocalDateTime debut, LocalDateTime fin);

    // Total des debits d'un compte sur une periode (0 s'il n'y en a aucun)
    @Query("""
        SELECT SUM(o.montant) FROM Operation o
         WHERE o.compte.numero = :numero AND o.type = 'DEBIT'
           AND o.date >= :debut AND o.date < :fin
        """)
    BigDecimal totalDebits(String numero, LocalDateTime debut, LocalDateTime fin);
}
\end{lstlisting}

Le service \texttt{ReleveService} les appelle pour un mois donné, avec pour début le 1\textsuperscript{er} du mois à 0\,h et pour fin le 1\textsuperscript{er} du mois suivant à 0\,h. Le relevé de mars ne doit contenir \textbf{que} des opérations de mars.

\begin{exobox}{Requête dérivée et requête JPQL}
Dans \texttt{OperationRepositoryTest} :
\begin{enumerate}[leftmargin=1.4em]
  \item Testez \texttt{findByCompteNumeroAndDateBetweenOrderByDateDesc} : opérations de deux comptes différents, ordre du résultat, et surtout les \textbf{valeurs limites} de la période. Où faut-il placer des opérations pour vérifier la règle « que des opérations de mars » ? (Indice : regardez la documentation de \texttt{Between} dans Spring Data.)
  \item Testez \texttt{totalDebits} : plusieurs débits, un crédit qui ne doit pas compter, une opération hors période, et un compte \textbf{sans aucun débit}.
  \item Ces deux méthodes contiennent chacune un bug. Pour chacun : notez dans \texttt{docs/tp4.md} le test qui le révèle et la cause, corrigez (dans la requête ou dans le nom de la méthode), relancez.
  \item Pourquoi ces bugs n'étaient-ils pas visibles dans les tests du TP3, où le repository était simulé ?
\end{enumerate}
\end{exobox}

\begin{exobox}{Isolation des tests}
\begin{enumerate}[leftmargin=1.4em]
  \item Écrivez un jeu de données \texttt{src/test/resources/donnees/comptes-et-operations.sql} (trois comptes, quelques opérations en février et en mars) et utilisez-le avec \texttt{@Sql} dans au moins deux tests.
  \item Écrivez deux tests qui créent chacun un compte \texttt{FR76-001}. Ils passent : pourquoi, alors que le numéro est \texttt{UNIQUE} ?
  \item Ajoutez \texttt{@Commit} sur l'un des deux et relancez la classe plusieurs fois. Que se passe-t-il ? Le résultat dépend-il de l'ordre ? Retirez \texttt{@Commit}.
\end{enumerate}
\end{exobox}

% ================================================================
\section{H2 n'est pas PostgreSQL (10 min)}

La migration \texttt{V3\_\_journal\_operations.sql} crée un \textbf{trigger} PL/pgSQL, dans l'esprit de ceux de R3.07 : chaque insertion dans \texttt{operation} ajoute une ligne dans la table \texttt{journal\_operation}.

\begin{exobox}{Le trigger, avec et sans PostgreSQL}
\begin{enumerate}[leftmargin=1.4em]
  \item Écrivez un test (\texttt{@DataJpaTest} + Testcontainers) qui enregistre une opération, force l'écriture en base avec \texttt{flush()} (via \texttt{TestEntityManager} ou \texttt{saveAndFlush}), puis vérifie avec un \texttt{JdbcTemplate} qu'une ligne est apparue dans \texttt{journal\_operation}. Pourquoi le \texttt{flush()} est-il nécessaire ?
  \item Copiez ce test dans une nouvelle classe, \textbf{sans} \texttt{Replace.NONE} ni conteneur : Spring utilise alors H2 (présent dans les dépendances de test). Lancez-le. Que se passe-t-il, et à quelle étape ?
  \item En deux ou trois phrases dans \texttt{docs/tp4.md} : quels autres éléments de votre base de R3.07 auraient posé le même problème ? Qu'en concluez-vous sur le choix de la base de test ? Supprimez ensuite cette classe.
\end{enumerate}
\end{exobox}

% ================================================================
\section{Tester l'API REST jusqu'à la base (25 min)}

L'application expose deux points d'accès :

\begin{center}
\renewcommand{\arraystretch}{1.25}
\small
\begin{tabular}{@{}l l >{\raggedright\arraybackslash}p{7.2cm}@{}}
\toprule
\texttt{GET} & \texttt{/api/comptes/\{numero\}/solde?devise=USD} & \texttt{200} + \texttt{\{"numero", "solde", "devise"\}} ; \texttt{404} si le compte n'existe pas \\
\texttt{POST} & \texttt{/api/virements} & corps \texttt{\{"source", "cible", "montant"\}} ; \texttt{201} si accepté, \texttt{404} si un compte n'existe pas, \texttt{422} si le solde est insuffisant, \texttt{400} si le montant est invalide \\
\bottomrule
\end{tabular}
\end{center}

\begin{infobox}[L'API est protégée]
Sans utilisateur connecté, l'API répond \texttt{401} ; un \texttt{POST} sans jeton CSRF reçoit \texttt{403}. Dans vos tests MockMvc, simulez un utilisateur avec \texttt{@WithMockUser("alice")} (sur la classe ou sur un test), et ajoutez \texttt{.with(csrf())} à chaque requête \texttt{POST} (imports : \texttt{org.springframework.security.test...}).
\end{infobox}

\begin{exobox}{MockMvc, PostgreSQL et doublures}
Créez \texttt{src/test/java/fr/iutvannes/banque/web/VirementApiTest.java} avec \texttt{@SpringBootTest}, \texttt{@AutoConfigureMockMvc} et le conteneur PostgreSQL. Remplacez par \texttt{@MockitoBean} \textbf{uniquement} ce qui sort de l'application : le service de change et le service de notification.
\begin{enumerate}[leftmargin=1.4em]
  \item \texttt{GET} solde : en euros, en dollars (taux simulé), compte inconnu.
  \item \texttt{POST} virement accepté : vérifiez le statut \textbf{et} les soldes \textbf{dans la base} après l'appel (rechargez les comptes depuis le repository).
  \item \texttt{POST} virement refusé : un cas pour chacun des statuts \texttt{404}, \texttt{422} et \texttt{400}. Les soldes en base doivent être inchangés.
  \item Un virement de plus de 1\,000\,\euro{} : la notification est-elle envoyée ? Comment le vérifier ici ?
  \item Activez un profil Spring \texttt{test} (\texttt{@ActiveProfiles("test")}) avec un fichier \texttt{src/test/resources/application-test.properties} qui affiche le SQL exécuté (\texttt{spring.jpa.show-sql=true}). Combien de requêtes SQL pour un virement ?
\end{enumerate}
\end{exobox}

% ================================================================
\section{Un virement à moitié fait ? (15 min)}

Un virement vers un compte \textbf{clôturé} doit être refusé : la méthode \texttt{crediter} de \texttt{Compte} lève une \texttt{CompteClotureException} (règle R4 du TP1), et l'API renvoie \texttt{422}.

\begin{exobox}{Vérifier la transaction}
\begin{enumerate}[leftmargin=1.4em]
  \item Écrivez le test : un compte source à 500\,\euro{}, un compte cible clôturé, \texttt{POST /api/virements} de 100\,\euro{}. Vérifiez le statut, puis le solde de la source \textbf{relu dans la base}.
  \item Si votre classe de test est annotée \texttt{@Transactional}, retirez cette annotation pour ce test (nettoyez la base avec un \texttt{@AfterEach} à la place). Relancez. Le résultat change-t-il ? Expliquez pourquoi un test \texttt{@Transactional} pouvait \textbf{cacher} le problème.
  \item La source a-t-elle été débitée alors que le virement a été refusé ? Si oui, regardez \texttt{VirementService.virer} : dans quel ordre les opérations sont-elles faites, et dans combien de transactions ? Corrigez (annotation \texttt{@Transactional} au bon endroit), puis vérifiez que le test passe.
\end{enumerate}
\end{exobox}

% ================================================================
\section{Rendu}

\begin{infobox}[À rendre avant \rdatelimite]
Sur le dépôt du binôme, un \textbf{tag} \texttt{tp4} pointant sur le commit rendu :
\begin{codebox}
git tag tp4
git push origin tp4
\end{codebox}
Le commit doit contenir :
\begin{itemize}[leftmargin=1.2em, itemsep=1pt]
  \item \texttt{CompteRepositoryTest}, \texttt{OperationRepositoryTest} (dont le trigger), \texttt{VirementApiTest} ;
  \item le jeu de données \texttt{comptes-et-operations.sql} et le profil \texttt{application-test.properties} ;
  \item les corrections de \texttt{OperationRepository} et de \texttt{VirementService} ;
  \item \texttt{docs/tp4.md} : temps mesurés, journal des bugs, H2 et PostgreSQL, explication de la transaction ;
  \item une suite complète qui passe : \texttt{./mvnw test} (Docker démarré).
\end{itemize}
\end{infobox}

\end{document}
